\section{Sformułowanie problemu}
\label{sec:problem}

Dla instrukcji naturalnojęzykowej $p$, która w naszym projekcie wygenerowana jest przez LLM - model Ollamy należy wygenerować sekwencję tokenów kodu $x_0$, która odpowiada zadaniu opisanemu przez instrukcję.
W czasie treningu znany kod referencyjny jest częściowo zastępowany tokenami \texttt{[MASK]}, a model otrzymuje zadanie przewidzenia oryginalnych tokenów w pozycjach objętych maską.

Najpierw tworzona jest uszkodzona sekwencję $x_t$ oraz maska $m$.
Maska (m) określa, które tokeny zostały ukryte: wartość (m\_i=1) oznacza pozycję zamaskowaną, natomiast (m\_i=0) — pozycję pozostawioną bez zmian.

Na podstawie uszkodzonej sekwencji model przewiduje rozkład prawdopodobieństwa tokenów dla każdej pozycji.
Funkcja straty jest jednak obliczana wyłącznie dla pozycji zamaskowanych, ponieważ to właśnie ich odtworzenia model ma się nauczyć:


\begin{equation}
\label{eq:ce}
\mathcal{L}_{CE}
= -\frac{1}{|M|}\sum_{i\in M}\log p_\theta(x_{0,i}\mid x_t,p,t),
\end{equation}

gdzie (M) jest zbiorem zamaskowanych pozycji, a (|M|) oznacza liczbę tych pozycji.
Wyrażenie z sumą oznacza prawdopodobieństwo, jakie model przypisał prawidłowemu tokenowi (x\_{0,i}) na podstawie uszkodzonej sekwencji (x\_t), kontekstu (p) oraz poziomu zaszumienia (t).
